% ============================================================
\section{Características de calidad}\label{sec:calidad}
% ============================================================

Esta sección convierte los stakeholders, concerns, restricciones y tensiones de las secciones~\ref{sec:stakeholders} a~\ref{sec:preguntas} en características de calidad. Cada característica se lleva a la subcaracterística de ISO/IEC~25010 que realmente importa para este juego, se justifica con los stakeholders y concerns que la exigen, se prioriza y se enlaza con los escenarios de la sección~\ref{sec:escenarios} que la vuelven verificable.

% ------------------------------------------------------------
\subsection{Criterio de selección}

Una característica entra en la lista solo si cumple las cuatro condiciones siguientes. Si falla alguna, se descarta o se absorbe en otra.

\begin{itemize}
  \item S1. Al menos un stakeholder la exige con un concern propio, dicho con sus palabras en la sección~\ref{sec:concerns}. No basta con que sea deseable en general.
  \item S2. Una restricción, dependencia o tensión la vuelve difícil de cumplir. Si cualquier arquitectura razonable la cumple sin esfuerzo, no es un driver.
  \item S3. Obliga a elegir entre alternativas de arquitectura: cambia qué componente decide, dónde vive la lógica o qué se persiste y cuándo.
  \item S4. Se puede expresar como un escenario con una medida comprobable.
\end{itemize}

% ------------------------------------------------------------
\subsection{Características seleccionadas}

\begin{tabla}{|L{1.05cm}|L{2.75cm}|L{3.3cm}|Y|L{1.55cm}|}
\hline
\filacab \cab{ID} & \cab{Característica (ISO/IEC 25010)} & \cab{Stakeholders y concerns} & \cab{Por qué es específica de este sistema} & \cab{Escenarios} \\ \hline
\endhead
QA-01 & Desempeño: \newline comportamiento temporal &
STK-01, STK-09. \newline CRN-02. \newline REQ-01, REQ-02, DEP-04. &
Es la única característica con un número impuesto: menos de 1 segundo (REQ-02). Colocar una barrera obliga a comprobar que ambos jugadores conservan un camino a su meta, y como el modo en línea es obligatorio, esa comprobación compite por el mismo segundo con el viaje por TCP y con la escritura en SQL Server (TEN-01, TEN-06). &
ESC-01 \\ \hline
QA-02 & Usabilidad: \newline protección contra errores de usuario y capacidad de aprendizaje &
STK-01, STK-09, STK-02. \newline CRN-01, CRN-03. \newline REQ-07. &
El criterio de aceptación no es un adulto, sino un niño de 8 años (REQ-07). La regla de no bloquear por completo el camino del rival es la menos intuitiva del juego y la que más jugadas inválidas produce, y el niño tiene que entender por qué se le rechaza sin leer un párrafo (CRN-01). &
ESC-02, ESC-03 \\ \hline
QA-03 & Seguridad: \newline integridad, confidencialidad y autenticidad &
STK-15, STK-05, STK-10. \newline CRN-10, CRN-11, CRN-22, CRN-23. \newline CON-02, CON-03, REQ-03. &
Tres restricciones del cliente la agravan a la vez: TCP sin cifrado propio (CON-02), protocolo hecho a la medida porque se prohíben los WebSockets (CON-03), lo que amplía la superficie de ataque, y un segundo factor que no sirve de nada si el código viaja en claro (REQ-03, Q-09). Además, el especialista asume que el cliente será modificado (CRN-23). &
ESC-03, ESC-05, ESC-06 \\ \hline
QA-04 & Flexibilidad: \newline adaptabilidad a otros idiomas y regiones (localizabilidad) &
STK-11, STK-08, STK-09. \newline CRN-15, CRN-16. \newline CON-05, REQ-06. &
El juego se distribuye internacionalmente y el cambio de idioma es obligatorio (REQ-06). Unicode tiene que funcionar de extremo a extremo (CON-05): pantallas, sockets y columnas de SQL Server. Choca con la interfaz para niños, que pide íconos con texto incrustado (TEN-05). &
ESC-11 \\ \hline
QA-05 & Mantenibilidad: \newline modificabilidad y modularidad &
STK-04, STK-05, STK-06, STK-07, STK-17, STK-18. \newline CRN-07, CRN-13, CRN-14, CRN-20, CRN-24, CRN-26. &
Hay cinco decisiones abiertas que caerán con el proyecto avanzado: IA (Q-01), LAN (Q-02), segundo factor (Q-03), nombre (Q-07) y distancia a Quoridor (Q-08). Con 14 semanas (CON-13), cada decisión tardía que obligue a rehacer el núcleo pone en riesgo la entrega (TEN-04). &
ESC-12 a ESC-16 \\ \hline
QA-06 & Fiabilidad: \newline recuperabilidad, tolerancia a fallos y disponibilidad &
STK-01, STK-13, STK-14. \newline CRN-06, CRN-19, CRN-21. \newline CON-04, CON-17, DEP-01, DEP-03, DEP-04. &
El modo en línea es obligatorio y depende de infraestructura externa (DEP-04). Los niños juegan desde redes domésticas inestables, y una caída del Wi-Fi no cierra la conexión TCP de forma ordenada (CON-17), así que ninguno de los dos extremos se entera por sí solo. La escritura asíncrona que protege la latencia abre una ventana de datos que se pueden perder (TEN-06). &
ESC-08, ESC-09, ESC-10 \\ \hline
QA-07 & Mantenibilidad: \newline capacidad de prueba &
STK-10, STK-16, STK-03. \newline CRN-08, CRN-12, CRN-17. \newline CON-14. &
Cada semana hay que entregar un incremento auditable (CON-14), y eso solo es viable si las partidas se ejecutan y se reproducen sin interfaz. Además, es la única forma de demostrar que cliente y servidor aplican la misma lógica (CRN-12). &
ESC-12, ESC-17 \\ \hline
QA-08 & Seguridad: \newline responsabilidad (accountability) y no repudio &
STK-01, STK-12, STK-13, STK-16, STK-17. \newline CRN-18, CRN-19. \newline REQ-05. &
La moderación por conteo de reportes (REQ-05) se puede abusar con reportes coordinados (TEN-02). Sin registro de quién hizo qué y cuándo, no se puede revisar una sanción (Q-11) ni saber qué se perdió en una caída (CRN-19). &
ESC-07, ESC-09 \\ \hline
QA-09 & Protección del usuario menor: \newline operación a prueba de fallos e identificación de riesgos (safety) &
STK-02, STK-17, STK-19, STK-01. \newline CRN-04, CRN-05, CRN-25. \newline REQ-04, REQ-05. &
Los usuarios son menores de edad y los padres pueden vetar el juego (CRN-04). El filtro de chat es obligatorio (REQ-04) y no puede evadirse, por lo que tiene que aplicarse en el servidor. A diferencia de QA-03, que protege al sistema, esta característica protege a la persona. &
ESC-04, ESC-07 \\ \hline
QA-10 & Flexibilidad: \newline instalabilidad \newline \textit{(condicionada a Q-10)} &
STK-19, STK-14. \newline CRN-25. \newline CON-16, DEP-08. &
Solo aplica si el juego se publica en una plataforma de distribución (Q-10). En ese caso, la plataforma rechaza el paquete si pide componentes de .NET (CON-16) que el equipo del usuario no tiene. &
ESC-18 \\ \hline
\end{tabla}

% ------------------------------------------------------------
\subsection{Prioridad}

Las características no se priorizan en bloque, porque dentro de una misma característica hay escenarios con consecuencias muy distintas para la arquitectura. La prioridad se asigna a cada escenario en el Utility Tree de la sección~\ref{sec:utility}, con dos criterios: la importancia para los stakeholders y la sensibilidad arquitectónica, es decir, cuánto depende el cumplimiento del escenario de las decisiones estructurales.

% ------------------------------------------------------------
\subsection{Características descartadas}

\begin{tabla}{|L{3.95cm}|Y|}
\hline
\filacab \cab{Característica} & \cab{Motivo del descarte} \\ \hline
\endhead
Escalabilidad (flexibilidad) &
Falla S1. Ningún stakeholder expresa un concern sobre número de partidas simultáneas ni sobre crecimiento de usuarios. Si Q-10 se resuelve a favor de publicar, habrá que reconsiderarla. \\ \hline
Portabilidad multiplataforma &
Falla S2. CON-01 y CON-16 fijan la plataforma; no es un objetivo del proyecto. La parte que sí importa, que el juego se instale en la plataforma elegida, se conserva como QA-10. \\ \hline
Eficiencia: utilización de recursos y capacidad &
Falla S2. Un tablero de 9 × 9 con 20 barreras consume memoria y procesador despreciables. El único recurso crítico es el tiempo, y ya está cubierto por QA-01. \\ \hline
Reusabilidad &
Falla S1. Ningún stakeholder está interesado en reutilizar componentes en otro producto. \\ \hline
Compatibilidad: interoperabilidad y coexistencia &
Falla S3. El sistema solo se integra con los servicios ya cubiertos en DEP, y esas integraciones se tratan como fallas de terceros en QA-06. \\ \hline
Adecuación funcional &
Falla S3. Que las reglas sean correctas es obligatorio en cualquier arquitectura, así que no sirve para elegir entre ellas. Lo que sí distingue arquitecturas, que cliente y servidor apliquen exactamente la misma lógica (CRN-12), queda cubierto por QA-05 y QA-07. \\ \hline
\end{tabla}

% ------------------------------------------------------------
\subsection{Tensiones entre características}

Las características seleccionadas no se pueden maximizar a la vez. Cada tensión de la sección~\ref{sec:tensiones} se traduce en un par de características que compiten, y los escenarios de la sección~\ref{sec:escenarios} fijan hasta dónde se puede ceder en cada una.

\begin{tabla}{|L{2.5cm}|L{1.8cm}|Y|L{2.2cm}|}
\hline
\filacab \cab{Características} & \cab{Tensión} & \cab{En qué compiten} & \cab{Escenarios que fijan el límite} \\ \hline
\endhead
QA-01 contra QA-03 & TEN-01 & Validar en el servidor y cifrar el tráfico protege la partida, pero cada paso consume parte del segundo disponible. & ESC-01, ESC-05, ESC-06 \\ \hline
QA-01 contra QA-06 y QA-08 & TEN-06 & Persistir cada jugada da recuperación y evidencia, pero cada escritura en SQL Server cuesta tiempo de respuesta. & ESC-01, ESC-09 \\ \hline
QA-02 contra QA-03 & TEN-03 & El segundo factor añade pasos y dispositivos que un niño de 8 años difícilmente administra. & ESC-03, ESC-16 \\ \hline
QA-02 contra QA-04 & TEN-05 & Lo más claro para un niño es el texto dentro del ícono, que es justamente lo que no se puede traducir. & ESC-02, ESC-11 \\ \hline
QA-05 contra el plazo (CON-13) & TEN-04 & Cada punto de extensión cuesta tiempo hoy a cambio de flexibilidad futura. & ESC-13, ESC-14 \\ \hline
QA-09 contra QA-08 & TEN-02 & Revisar sanciones con justicia exige conservar mensajes de menores, lo que la protección de sus datos quiere evitar. & ESC-04, ESC-07 \\ \hline
\end{tabla}

% ------------------------------------------------------------
\subsection{Alcance de algunas características}

\begin{itemize}
  \item QA-08 se entiende como responsabilidad (accountability) del producto y no como auditabilidad del proceso. La parte de proceso, que cada semana haya un incremento demostrable (CRN-08, CON-14), se cubre con QA-05 y QA-07. La trazabilidad de decisiones (CRN-09) se cumple con la documentación y no con el producto.
  \item QA-09 se trata como protección (safety): lo que importa es que el chat falle de forma segura cuando algo no funciona.
  \item QA-10, instalabilidad, se incluye condicionada a Q-10, porque CRN-25 impone requisitos de empaquetado que ninguna otra característica cubre.
\end{itemize}
